%% LyX 2.4.4 created this file.  For more info, see https://www.lyx.org/.
%% Do not edit unless you really know what you are doing.
\documentclass[english]{article}
\usepackage[T1]{fontenc}
\usepackage[latin9]{inputenc}
\usepackage{units}
\usepackage{enumitem}
\usepackage{amsmath}
\usepackage{amsthm}
\usepackage{geometry}
\geometry{verbose,tmargin=1in,bmargin=1in,lmargin=1in,rmargin=1in}

\makeatletter
%%%%%%%%%%%%%%%%%%%%%%%%%%%%%% Textclass specific LaTeX commands.
\numberwithin{equation}{section}
\numberwithin{figure}{section}
\newlength{\lyxlabelwidth}      % auxiliary length 

%%%%%%%%%%%%%%%%%%%%%%%%%%%%%% User specified LaTeX commands.
\usepackage{fancyhdr}
\usepackage{lastpage}
\pagestyle{fancy}
\usepackage{enumitem}
\usepackage{ifthen}
\everymath=\expandafter{\the\everymath\displaystyle}
%\DeclareMathSizes{10}{10}{10}{10}
\usepackage{multicol}

\usepackage{tikz}
\usepackage{tikz-3dplot}
\usetikzlibrary{calc,arrows}

\usetikzlibrary{patterns}
\usetikzlibrary{plotmarks}
\usepackage{pgfplots}
\pgfplotsset{compat=newest}

\usepackage{calc}
\usepackage[nomessages]{fp}

\usepackage{datetime}
% Added by lyx2lyx
\setlength{\parskip}{\smallskipamount}
\setlength{\parindent}{0pt}

\makeatother

\usepackage{babel}
\begin{document}
\renewcommand{\author}{Dr.\,James Ashe}

\newcommand{\classNum}{331}

\newcommand{\classSec}{A}

\newcommand{\examType}{Exam 1 Review: 9.1-9.4, 10.1}

\renewcommand{\date}{\formatdate{9}{9}{2013}}

%%First page's header%%%%%%%%%%%%%%%%%%%%%%
\fancypagestyle{firststyle} %%header settings for first pagr
{
	\fancyhead[L]{MTH \classNum{}(\classSec{}) \author{}\\
	\textbf{\examType{}} ~\date{}}
	\fancyhead[C]{}
	\fancyhead[R]{\textbf{Name:\hspace{4cm}}}
	\setlength{\headsep}{14pt}
}
%%%%%%%%%%%%%%%%%%%%%%%%%%%%%%%%%%%%%%%%%%

%%Next page's header%%%%%%%%%%%%%%%%%%%%%%%
\setlist{nolistsep}
\lhead{\classNum{}(\classSec{})} 
\chead{\textbf{Name:}} 
\cfoot{\thepage{}~of~\pageref{LastPage}} 
\setlength{\headsep}{5pt} 
%%%%%%%%%%%%%%%%%%%%%%%%%%%%%%%%%%%%%%%%%%%

%%Various settings; see the preamble for other settings%%%%%
%\setenumerate[0]{leftmargin=12pt,itemindent=0pt,labelwidth=0pt}
\setlist{topsep=.5pt}
\renewcommand{\labelenumi}{\textbf{\arabic{enumi}.}}
\renewcommand{\labelenumii}{\textbf{\alph{enumii}.}}
\thispagestyle{firststyle}
%%%%%%%%%%%%%%%%%%%%%%%%%%%%%%%%%%%%%%%%%%

\newcommand{\xmax}{0}
\newcommand{\xmin}{-\xmax}
\newcommand{\xstep}{0}
\newcommand{\ymax}{0}
\newcommand{\ymin}{-\ymax}
\newcommand{\ystep}{0} 
\newcommand{\dspace}{\vspace{1cm}}
\newcommand{\xa}{0}
\newcommand{\xb}{0}
\newcommand{\ya}{1}
\newcommand{\yb}{1}

\textbf{Justify your answers and show your work. All questions are
of equal value.}
\begin{enumerate}[resume]
\item Write the infinite series, $0.3+0.03+0.003+\cdots$, using sigma notation
(i.e. using $\Sigma$) . 

\begin{enumerate}[resume]
\item []${\displaystyle \sum_{k=1}^{\infty}\frac{3}{10^{k}}}$
\end{enumerate}
\item Write the first 3 terms of the sequence of partial sums of the series
${\displaystyle \sum_{k=1}^{\infty}\frac{3^{k}}{2}}$.

\begin{enumerate}[resume]
\item []$\sum_{k=1}^{3}\frac{3^{k}}{2}=\frac{3}{2}+\frac{9}{2}+\frac{27}{2}$\vfill
\end{enumerate}
\end{enumerate}
\emph{Evaluate the geometric series or state that it diverges. }
\begin{enumerate}[resume]
\item ${\displaystyle \sum_{k=0}^{\infty}\frac{2^{k}}{7^{k}}}$.

\begin{enumerate}[resume]
\item []${\displaystyle \sum_{k=0}^{\infty}\frac{2^{k}}{7^{k}}=\sum_{k=0}^{\infty}\left(\frac{2}{7}\right)^{k}=\frac{1}{1-\nicefrac{2}{7}}=\frac{7}{5}}$
since $r=\frac{2}{7}$ and $a=1$. 
\end{enumerate}
\item ${\displaystyle \sum_{m=2}^{\infty}\left(\frac{1}{3}\right)^{m}}$.

\begin{enumerate}[resume]
\item []${\displaystyle \sum_{m=2}^{\infty}\left(\frac{1}{3}\right)^{m}={\displaystyle \sum_{m=0}^{\infty}\left(\frac{1}{3}\right)^{m+2}={\displaystyle \sum_{m=2}^{\infty}\left(\frac{1}{3}\right)^{2}\left(\frac{1}{3}\right)^{m}}}=\frac{\nicefrac{1}{9}}{1-\nicefrac{1}{3}}=\frac{1}{6}}$
since $r=\frac{1}{3}$ and $a=\frac{1}{9}$. Note that you can also
just write out a few terms to determine the $a$ and $r$ . \vfill
\end{enumerate}
\end{enumerate}
\emph{Find the limit of the following sequences or determine that
the limit does not exist. }
\begin{enumerate}[resume]
\item ${\displaystyle \lim_{n\rightarrow\infty}\left\{ \frac{n^{4}-1}{n^{4}+10}\right\} }$.

\begin{enumerate}[resume]
\item []${\displaystyle \frac{n^{4}-1}{n^{4}+10}}\frac{\nicefrac{1}{n^{4}}}{\nicefrac{1}{n^{4}}}=\frac{1-\nicefrac{1}{n^{4}}}{1+\nicefrac{10}{n^{4}}}$$\Rightarrow{\displaystyle \lim_{n\rightarrow\infty}\left\{ \frac{n^{4}-1}{n^{4}+10}\right\} }={\displaystyle \lim_{n\rightarrow\infty}\left\{ \frac{1-\nicefrac{1}{n^{4}}}{1+\nicefrac{10}{n^{4}}}\right\} =1}$
\end{enumerate}
\item ${\displaystyle \lim_{n\rightarrow\infty}\left\{ \frac{n^{0.01}}{\ln n}\right\} }$.

\begin{enumerate}[resume]
\item []$\ln n\ll n^{0.01}$ so then $\lim_{n\rightarrow\infty}\left\{ \frac{n^{0.01}}{\ln n}\right\} =\infty$. \vfill
\end{enumerate}
\end{enumerate}
\emph{Use the Divergence Test to determine whether the following series
diverge, or state that the test is inconclusive (use your results
from above).}
\begin{enumerate}[resume]
\item ${\displaystyle \sum_{k=1}^{\infty}\frac{k^{4}-1}{k^{4}+10}}$.

\begin{enumerate}[resume]
\item []From above, ${\displaystyle \lim_{k\rightarrow\infty}}\left\{ \frac{k^{4}-1}{k^{4}+10}\right\} \neq0$,
so by the Divergence theorem ${\displaystyle \sum_{k=1}^{\infty}\frac{k^{4}-1}{k^{4}+10}}$
diverges.
\end{enumerate}
\item ${\displaystyle \sum_{k=1}^{\infty}\frac{k^{0.01}}{\ln k}}$.

\begin{enumerate}[resume]
\item []From above $\lim_{n\rightarrow\infty}\left\{ \frac{n^{0.01}}{\ln n}\right\} \neq0$
so by the Divergence theorem ${\displaystyle \sum_{k=1}^{\infty}\frac{k^{0.01}}{\ln k}}$
diverges.
\end{enumerate}
\end{enumerate}
\emph{Using the Squeeze Theorem, find the limit of the following sequence
or determine that the limit does not exist.}
\begin{enumerate}[resume]
\item ${\displaystyle \lim_{n\rightarrow\infty}\left\{ \frac{\sin^{2}n}{2^{n}}\right\} }$.

\begin{enumerate}[resume]
\item []$-1\leq\sin x\leq1\Rightarrow0\leq\sin^{2}x\leq1$. So then $0\leq\frac{\sin^{2}n}{2^{n}}\leq\frac{1}{2^{n}}\Rightarrow{\displaystyle 0=\lim_{n\rightarrow\infty}}\left\{ 0\right\} \leq{\displaystyle \lim_{n\rightarrow\infty}}\left\{ \frac{\sin^{2}n}{2^{n}}\right\} \leq{\displaystyle \lim_{n\rightarrow\infty}}\left\{ \frac{1}{2^{n}}\right\} =0$.
So then ${\displaystyle \lim_{n\rightarrow\infty}\left\{ \frac{\sin^{2}n}{2^{n}}\right\} }=0$
by the Squeeze theorem. \vfill
\end{enumerate}
\end{enumerate}
\emph{Use the Integral Test to determine whether the series converge,
diverge, or state why the test cannot be used. Check that the conditions
of the test are satisfied.}
\begin{enumerate}[resume]
\item ${\displaystyle \sum_{k=1}^{\infty}\frac{k}{e^{k}}}$. For $f(x)=\frac{x}{e^{x}}$,
$f'(x)=\left(1-x\right)e^{-1}$ and $\int f(x)\,dx=2xe^{-1}$.

\begin{enumerate}[resume]
\item []For $x\geq1$, $f'(x)=\left(1-x\right)e^{-1}\leq0$, $f(x)=\frac{x}{e^{x}}\geq0$,
and $f(x)$ is continuous. So by the integral test it diverges since,
\[
{\displaystyle \int_{1}^{\infty}\frac{x}{e^{x}}=\lim_{b\rightarrow\infty}\left(2xe^{-1}\vert_{1}^{b}\right)={\displaystyle \lim_{b\rightarrow\infty}2xe^{-1}-2e^{-1}=\infty-2e^{-1}=\infty\mbox{.}}}
\]
\end{enumerate}
\item ${\displaystyle \sum_{k=2}^{\infty}\frac{1}{(2x-3)^{2}}}$. For $f(x)=\frac{1}{(2x-3)^{2}}$,
$f'(x)=\frac{-4}{(2x-3)^{2}}$ and $\int f(x)\,dx=\frac{-1}{2(2x-3)}$.

\begin{enumerate}[resume]
\item []For $x\geq2$, $f(x)$ is continuous, $f'(x)\leq0$, and $f(x)\geq0$.
So by the integral test the series converges since 
\[
{\displaystyle \int_{2}^{\infty}\frac{1}{(2x-5)^{2}}\,dx={\displaystyle \lim_{b\rightarrow\infty}\frac{1}{2(2x-3)}|_{2}^{b}=0-(\nicefrac{-1}{2})=\nicefrac{1}{2}}}
\]
\vfill
\end{enumerate}
\item Let $f(x)=2\cos(x)$.

\begin{enumerate}[resume]
\item Find the Taylor polynomial of order $n=3$ for $f(x)$ centered at
$0$.

\begin{enumerate}[resume]
\item []$f'(x)=-2\sin(0)=0$, $f''(x)=-2\cos(0)=-2$, $f^{(3)}(0)=2\sin(x)=0$.
So then $p_{3}(x)=2+0-x^{2}$. 
\end{enumerate}
\item Use the Taylor polynomial from above to approximate $2\cos(0.1)$.

\begin{enumerate}[resume]
\item []$p_{3}(0.1)=2-(.1)^{2}=1.99$
\end{enumerate}
\item Compute the absolute error in the approximation assuming the exact
value is given by a calculator.

\begin{enumerate}[resume]
\item []$\left|p_{3}(0.1)-2\cos(0.1)\right|\approx.00001$
\end{enumerate}
\end{enumerate}
\end{enumerate}

\end{document}
